% =====================================================================
%  CHAPTER 1 — INTRODUCTION
%  Every number in this chapter comes with its source and with the
%  experimental environment that produced it. No measurement of this
%  thesis appears in this chapter.
%
%  ENGLISH IS THE MASTER. content/el/ch1.tex follows this file, block
%  for block. See docs/SYNC.md.
% =====================================================================

%%blk intro.head
\chapter{Introduction: the consensus problem and the question of this thesis}
\label{ch:intro}

%%blk intro.head.roadmap
This chapter does four things, in this order. It states the problem that
consensus solves and the two obstacles that make it hard, faults and
delays. It explains why one protocol never settled the question, and why
the published comparisons between protocols do not transfer from one
paper to the next. It then states the aim of this thesis, its research
questions and its limits. It closes with the rule this thesis follows
about the quality of its sources, because a comparison is worth exactly
as much as the evidence under it.

%%blk intro.problem.head
\section{Agreement under faults and delays}
\label{sec:problem}

%%blk intro.problem.01
A distributed system consists of processes that are separated in space
and that communicate only by messages. A system is distributed when the
transmission delay is not negligible against the time between two events
of the same process \cite{lamport1978time}. This definition has one
consequence: no global clock exists. For two events in different
processes it is sometimes impossible to say which one happened first.
The only inherent order relation is the partial \emph{happened-before}
relation \cite{lamport1978time}. A service that must behave like one
single reliable machine must therefore build that order itself, because
the network does not supply it.

%%blk intro.problem.02
The standard abstraction that builds this order is \emph{state machine
replication} (SMR). SMR models the service as a deterministic state
machine and replicates it on $n$ nodes. If all correct nodes execute the
same sequence of commands in the same order, they reach the same state
and produce the same outputs \cite{schneider1990implementing}. The
reliability problem then reduces to one single coordination problem,
\emph{consensus}: the agreement of the correct nodes on one common value
for each position of the sequence. The correctness criterion that the
client expects from the replicated service is \emph{linearizability}.
Every operation appears to happen at one instant between its invocation
and its response, and operations that do not overlap in time keep their
order \cite{herlihy1990linearizability}.

%%blk intro.problem.03
Two properties judge a consensus protocol. \emph{Safety} requires that
two correct nodes never decide differently for the same position.
\emph{Liveness} requires that the system decides every request of a
correct client eventually. The fundamental obstacle is that no
deterministic protocol holds both properties at the same time in a fully
asynchronous system. Fischer, Lynch and Paterson proved that
\emph{every deterministic} protocol for this problem can fail to
terminate, even with one single faulty process
\cite{fischer1985impossibility}. The result is
structural and not a technical limitation: in an asynchronous system no
process can tell a dead process from a slow one.

%%blk intro.problem.04
All the protocols that this thesis examines take the same way out, the
model of \emph{partial synchrony} \cite{dwork1988consensus}. In the form
that modern protocols adopt, an unknown moment exists. It is the GST, the
\emph{Global Stabilization Time}. After the GST, the network delivers
messages between correct nodes within a known bound $\Delta$. The
Byzantine protocols of this thesis state the model in exactly this form
\cite{yin2019hotstuff} and \cite{malkhi2023hotstuff2}\prep. Safety holds
at all times, even before the GST, and liveness is guaranteed only after
it.

%%blk intro.problem.05
The consequence for every experimental evaluation is important. The
recovery time after the loss of the leader depends to a large degree on
the timer settings, not on the algorithm alone. Two works state this
explicitly. The reproducibility study of Raft notes it first
\cite{howard2015raftrefloated}. Berger et al.\ repeat it for Byzantine
systems. They write that the recovery time depends on the timer
parameters, and that they \emph{do not exclude} that tighter settings
would be feasible \cite{berger2023simulation}.

%%blk intro.problem.06
The second dimension that organizes the whole literature is the
\emph{fault model}. In the \emph{crash-fault} model a node can only
stop. In the Byzantine model a node can behave in an arbitrary way, send
contradictory messages to different receivers, and coordinate with other
faulty nodes \cite{lamport1982byzantine}. The distinction is not
philosophical. It fixes how many nodes a system needs, as
Section~\ref{sec:notonce} shows.

%%blk intro.notonce.head
\section{Why the problem was not solved once}
\label{sec:notonce}

%%blk intro.notonce.01
Viewstamped Replication appeared in 1988. It is the first protocol of the
line that this thesis examines. Almost four decades have passed since
then, and the production of new protocols has not stopped.
Figure~\ref{fig:timeline} gives the order in time under one explicit
rule: each protocol sits at its \emph{first refereed publication}, in
conference proceedings or in a journal. Three of the ten protocols have
no such publication. The figure draws them with an empty marker, puts
them at the first public release of their design, and marks their source.
This thesis states the rule here, because the word
\enquote{chronologically} without a rule becomes arbitrary, and the
arbitrariness shows.

%%blk intro.notonce.fig.timeline
\dgfigure{fig:timeline}%
{Timeline of the protocols that this thesis examines, placed by the rule
of first refereed publication.}%
{Timeline of the protocols that this thesis examines, placed by the rule
of first refereed publication. \textbf{Filled marker} = the first
refereed publication, which also fixes the position: VR
\cite{oki1988viewstamped}, Paxos \cite{lamport1998parttime}, PBFT
\cite{castro1999pbft}, ZAB \cite{junqueira2011zab}, EPaxos
\cite{moraru2013epaxos}, Raft \cite{ongaro2014raft}, HotStuff
\cite{yin2019hotstuff}. \textbf{Empty marker} = no refereed publication
exists. The position then comes from the first public release of the
design, and the figure marks it. KRaft: KIP-500, technical
documentation\grey\ of 2019 \cite{kafka2019kip500}. General
availability came later, with Kafka~3.3; the date of that release is not
documented by the source cited here, and it does not fix the position in
the figure. Tendermint: algorithmic preprint\prep\ of
2018 \cite{buchman2018gossip}, whitepaper\grey\ of 2014
\cite{kwon2014tendermint}. A refereed paper does exist
\cite{cason2021tendermint}, but it appeared in 2021 and describes the
design, architecture and performance of the \emph{network}; placing the
protocol there would put it seven years after the design it implements.
The figure therefore places Tendermint at its algorithmic preprint and
marks the marker as empty, and this paragraph is the reservation that
goes with that choice. HotStuff-2: extended
abstract on ePrint\prep\ \cite{malkhi2023hotstuff2}. \emph{Two
reservations about provenance.} The work \cite{lamport1998parttime}
carries a received date of January 1990 and an accepted date of March
1998, but the two copies of the corpus \emph{word it} differently (see
\S\ref{sec:sources}). The technical report SRC~RR-49 is widely dated to
1989. We did not locate it as evidence, and for that reason it does not
fix the position. PBFT also has an extended journal version in 2002
\cite{castro2002pbftrecovery}. Theoretical milestones:
\cite{lamport1978time,lamport1982byzantine,fischer1985impossibility,%
dwork1988consensus,schneider1990implementing,herlihy1990linearizability}.}%
{% fig-timeline — first-refereed-publication rule, CFT above the axis, BFT below.
% Structure figure: dgflow. The only geometry it sets is the year-to-length
% mapping, and that is tied to \dgwidth so the timeline is exactly as wide as
% every other figure in the thesis: 51 years (1977-2028) across \dgwidth.
% Every human-readable string is a label key; protocol names, venues and years
% are identical in both editions and stay literal.
\begin{dgflow}[x={\dgwidth/51}]
  \def\yzero{1977}

  % Theoretical milestones. A third category — neither CFT nor BFT protocols —
  % so dgAux rather than either series colour.
  \node[dgnote, text=dgAux, anchor=west] at (0,4.45)
    {\textbf{\fl{timeline.milestones.title}}\\
     1978 \fl{timeline.milestone.clocks}\, ·\, 1982 \fl{timeline.milestone.byzmodel}\, ·\, 1985 FLP\\
     1988 \fl{timeline.milestone.partialsync}\, ·\, 1990 SMR\, ·\, 1990 \fl{timeline.milestone.linearizability}};

  % The time axis and its five-year ticks. The year LABELS are drawn last, so
  % that their dgtag background sits over the leader lines rather than under.
  \draw[dgaxis] (0,0) -- (51,0);
  \foreach \yr in {1980,1985,1990,1995,2000,2005,2010,2015,2020,2025}{
    \pgfmathsetmacro{\xx}{\yr-\yzero}
    \draw[dgaxis,-] (\xx,-0.08) -- (\xx,0.08);}

  % Which half is which. At the left end, where the axis is still empty.
  \node[dgtag, fill=dgPaper, text=dgCFT, anchor=west] at (0.4, 1.00) {\textbf{CFT}};
  \node[dgtag, fill=dgPaper, text=dgBFT, anchor=west] at (0.4,-1.00) {\textbf{BFT}};

  % Crash-fault protocols, above the axis. Three stagger levels: 2011-2019 is
  % 1.7 cm of axis carrying 5.6 cm of label, and the anchors are chosen so that
  % no leader line ever runs through another protocol's label.
  % year / level / anchor / filled marker? / label
  \foreach \yr/\lvl/\anc/\solid/\lbl in {%
      1988/1.00/{south}/{1}/{VR\\PODC '88},%
      1998/1.00/{south}/{1}/{Paxos\\TOCS '98},%
      2011/1.00/{south east}/{1}/{ZAB\\DSN '11},%
      2013/1.95/{south east}/{1}/{EPaxos\\SOSP '13},%
      2014/2.90/{south}/{1}/{Raft\\ATC '14},%
      2019/1.00/{south west}/{0}/{KRaft\\KIP-500 '19\grey}}{
    \pgfmathsetmacro{\xx}{\yr-\yzero}
    \draw[dgarrow,-,draw=dgCFT] (\xx,0.07) -- (\xx,\lvl);
    \ifnum\solid=1 \fill[dgCFT] (\xx,0.07) circle (1.7pt);
    \else \draw[dgarrow,-,draw=dgCFT,fill=dgPaper] (\xx,0.07) circle (1.7pt);\fi
    \node[dgtag, fill=dgPaper, text=dgCFT, anchor=\anc] at (\xx,\lvl) {\lbl};}

  % Byzantine-fault protocols, below the axis. Same three levels, same reason:
  % 2018, 2019 and 2023 are 1.1 cm apart and their labels are 1.5 cm wide.
  \foreach \yr/\lvl/\anc/\solid/\lbl in {%
      1999/-1.00/{north}/{1}/{PBFT\\OSDI '99},%
      2018/-1.00/{north east}/{0}/{Tendermint\\arXiv '18\prep},%
      2019/-1.95/{north east}/{1}/{HotStuff\\PODC '19},%
      2023/-2.90/{north}/{0}/{HotStuff-2\\ePrint '23\prep}}{
    \pgfmathsetmacro{\xx}{\yr-\yzero}
    \draw[dgarrow,-,draw=dgBFT] (\xx,-0.07) -- (\xx,\lvl);
    \ifnum\solid=1 \fill[dgBFT] (\xx,-0.07) circle (1.7pt);
    \else \draw[dgarrow,-,draw=dgBFT,fill=dgPaper] (\xx,-0.07) circle (1.7pt);\fi
    \node[dgtag, fill=dgPaper, text=dgBFT, anchor=\anc] at (\xx,\lvl) {\lbl};}

  % Year labels last: 1988, 1999, 2011, 2014 and 2019 all sit within half a
  % label of a five-year tick, so something has to be on top and it should be
  % the year.
  \foreach \yr in {1980,1985,1990,1995,2000,2005,2010,2015,2020,2025}{
    \pgfmathsetmacro{\xx}{\yr-\yzero}
    \node[dgtag, fill=dgPaper, text=dgGrey, anchor=north] at (\xx,-0.16) {\yr};}
\end{dgflow}
}

%%blk intro.notonce.02
The repetition does not happen because the earlier solutions failed to be
correct. It happens because correctness is only one dimension. Every
protocol is a \emph{position} in a design space, and every improvement
there is bought at a specific price. The recent unifying work of Amiri
et al.\ states this explicitly. It describes Byzantine protocols as
points of a space with more than ten design dimensions. In that space
each choice is a one-to-one function from protocol to protocol
\cite{amiri2024bedrock}.

%%blk intro.notonce.03
The clearest sign of this is the arithmetic of the quorums
(Figure~\ref{fig:quorum}). The node count does not rise arbitrarily. It
moves only when it buys something specific that can be named. With crash
faults a majority quorum is enough. A cluster of $n = 2f+1$ nodes stays
available while a majority works, so it tolerates $f$ faults
\cite{ongaro2014raft}. The general statement of safety as the
intersection of majority sets goes back to \cite{lamport1998parttime}.

%%blk intro.notonce.04
With Byzantine faults and messages that are not authenticated, a system
needs at least $n = 3f+1$ nodes. The bound appears first in 1980, as the
problem of interactive consistency \cite{pease1980reaching}. The allegory
of the Byzantine generals came later \cite{lamport1982byzantine}. The
same bound holds for the partially synchronous protocols examined here
\cite{dwork1988consensus,amiri2024bedrock}. There, however, it does not
follow from the absence of signatures, and Section~\ref{sec:bounds} shows
why.

%%blk intro.notonce.05
Trusted hardware (a \emph{trusted execution environment}) stops a faulty
node from sending contradictory messages without detection. If such
hardware is added, the bound returns to $2f+1$ \cite{amiri2024bedrock}.
Saving one phase of communication cost $5f+1$ nodes at first
\cite{martin2006fab}\unver. With authentication, $5f-1$ nodes are enough. That
number is also the proven lower bound for two-step Byzantine consensus
\cite{kuznetsov2021optimal}\unver, and practice reaches it
\cite{amiri2024bedrock}. Finally, support for \emph{order-fairness} needs
at least $4f+1$ nodes in its \emph{least} demanding parameterization. As
the parameter $\gamma$ comes near $0.5$ and the fairness requirement
becomes stricter, the bound $n > 4f/(2\gamma-1)$ rises without limit
\cite{kelkar2020orderfairness,amiri2024bedrock}\unver.

%%blk intro.notonce.fig.quorum
\dgfigure{fig:quorum}%
{Quorum arithmetic: the required number of nodes $n$ as a function of the
tolerated faults $f$.}%
{Quorum arithmetic: the required number of nodes $n$ as a function of the
tolerated faults $f$. These are not measurements. They are pure
arithmetic from bounds that the literature states. For $f=3$: Byzantine
tolerance costs 10 nodes against 7, fair ordering costs 13, and one phase
fewer costs 16. \emph{Sources.} The work \cite{amiri2024bedrock} collects
all five bounds: $3f{+}1$, $2f{+}1$ and $5f{\pm}1$ in \S3.2 E1, trusted
hardware in App.~A, O7, and fair ordering in \S3.3 Q1 and in design
choice 13. For these bounds it is a secondary source. For $3f{+}1$ it
gives five references, of which this thesis names one,
\cite{lamport1982byzantine}; the other four are Bracha 1984, Bracha and
Toueg 1985, Correia et al.\ 2006 and \cite{dwork1988consensus}. The
primary statement is earlier still, from 1980 \cite{pease1980reaching},
which \cite{amiri2024bedrock} does \emph{not} cite --- that attribution
is this thesis's own, and Section~\ref{sec:bounds} gives the argument
for it. For the remaining bounds the survey attributes $5f{+}1$ to
\cite{martin2006fab}\unver, $5f{-}1$ to \cite{kuznetsov2021optimal}\unver\
and \cite{abraham2021goodcase}\unver\ (it cites both), and $4f{+}1$ to
\cite{kelkar2020orderfairness}\unver. Of these sources, only
\cite{pease1980reaching} and \cite{lamport1982byzantine} had their
bibliographic data verified, through an index. The other four were not
verified, and their data come from the reference list of
\cite{amiri2024bedrock}. \emph{None of the six is in the corpus.}
\emph{Note on precision.} The work \cite{amiri2024bedrock} separates the
size of the quorum from the size of the network. In designs with active
and passive nodes the quorum falls to $2f{+}1$ while $n$ stays
$3f{+}1$.}%
{% fig-quorum — required n as a function of tolerated faults f. Pure arithmetic
% from bounds stated in the literature; no measurements.
% Data figure: dgdata. The five series are plotted in the order that makes the
% standard cycle list hand them their meanings: dgCFT crash, dgBFT Byzantine,
% then dgAux / dgSand / dgPurple for the three further bounds.
\begin{dgdata}[
  xlabel={\fl{quorum.xlabel} $f$},
  ylabel={\fl{quorum.ylabel} $n$},
  xmin=0.7, xmax=5.3, ymin=0, ymax=28,
  xtick={1,2,3,4,5}, ytick={0,5,10,15,20,25},
  legend style={at={(0.5,-0.30)}, anchor=north}, legend columns=1,
]
\addplot coordinates {(1,3)(2,5)(3,7)(4,9)(5,11)};
  \addlegendentry{$2f{+}1$ — \fl{quorum.crash}}
\addplot coordinates {(1,4)(2,7)(3,10)(4,13)(5,16)};
  \addlegendentry{$3f{+}1$ — \fl{quorum.byzantine}}
\addplot coordinates {(1,5)(2,9)(3,13)(4,17)(5,21)};
  \addlegendentry{$4f{+}1$ — \fl{quorum.fairorder} ($\gamma=1$)}
\addplot coordinates {(1,6)(2,11)(3,16)(4,21)(5,26)};
  \addlegendentry{$5f{+}1$ — \fl{quorum.onephaseless} (FaB)}
\addplot coordinates {(1,4)(2,9)(3,14)(4,19)(5,24)};
  \addlegendentry{$5f{-}1$ — \fl{quorum.twostep}}
\end{dgdata}
}

%%blk intro.notonce.06
The same logic governs the second large dimension, the \emph{view
change}. On the fault-free path the protocols differ mainly in topology
and in the number of phases. They diverge much more in the replacement of
the leader, which the designers of HotStuff call \enquote{the epicenter
of replication} \cite{yin2019hotstuff}. Table~\ref{tab:viewchange}
summarizes four representative positions on this axis.

%%blk intro.notonce.tab.viewchange
\begin{table}[t]
\centering
\caption[Four representative positions on the view-change axis: cost,
responsiveness and number of phases.]%
{Four representative positions on the view-change axis: cost,
responsiveness and number of phases. The table reproduces Table~1 of
\cite{malkhi2023hotstuff2}\prep\ together with the reservation that the
same source states. It must not be read without that reservation. It is
not exhaustive, because the same source lists other designs on the same
axis.}
\label{tab:viewchange}
\small
\begin{tabular}{@{}l c c c@{}}
\toprule
\textbf{Protocol} & \textbf{View-change cost} &
\textbf{Responsive} & \textbf{Phases} \\
\midrule
PBFT \cite{castro1999pbft}        & $O(n^{2})$ & yes & 2 \\
Tendermint \cite{buchman2018gossip}\prep & $O(n)$ & no & 2 \\
HotStuff \cite{yin2019hotstuff}   & $O(n)$     & yes & 3 \\
HotStuff-2 \cite{malkhi2023hotstuff2}\prep & $O(n)$ & yes & 2 \\
\bottomrule
\end{tabular}

\vspace{4pt}
\begin{minipage}{0.94\linewidth}\setstretch{1.08}\footnotesize
\textbf{The reservation of the source itself.} The source
\cite{malkhi2023hotstuff2}\prep\ states explicitly that PBFT and
Tendermint have a view-change cost of $O(n^{3})$ and $O(n^{2})$
respectively, and that the bounds of the table follow after both are
improved by a factor of $O(n)$ --- an improvement that the same source
attributes elsewhere to signature aggregation. The detailed count of
complexity in \cite{yin2019hotstuff} (Table~1) gives PBFT $O(n^{2})$
authenticators with a stable leader, $O(n^{3})$ on a view change, and
$O(f n^{3})$ over a sequence of $f$ leader failures. This is not an
independent confirmation, because the two works share an author. The work
\cite{jalalzai2023fasthotstuff} returns to two rounds and keeps
responsiveness, but it states explicitly that it does \emph{not} avoid
the transfer of a quadratic number of view-change messages when the
leader fails. What it reduces is the number of signatures that each node
verifies, from $O(n)$ to two. Malkhi and Nayak publish an extended
abstract on ePrint that has not been refereed
\cite{malkhi2023hotstuff2}\prep. It is cited as such.

\smallskip
\textbf{A phase is not a message delay.} The table counts
\emph{phases}, with the convention of its source. One phase of HotStuff
costs two consecutive delays, because it passes through the leader. One
voting phase of PBFT costs one delay, because it is all-to-all. The two
units do not compare, and Section~\ref{sec:performance} defines the
distinction once.
\end{minipage}
\end{table}

%%blk intro.notonce.07
The table shows that none of the four protocols is \enquote{better}. Each
one pays a different bill. Tendermint removes the view-change
subprotocol and pays with a wait for the worst possible $\Delta$. That
is, it pays with \emph{responsiveness}, the property that lets a protocol
advance at the real network time and not at the maximum network time
\cite{buchman2018gossip}\prep\ and \cite{malkhi2023hotstuff2}\prep.
HotStuff pays with a third phase \cite{yin2019hotstuff}. HotStuff-2
claims that neither of the two prices was necessary
\cite{malkhi2023hotstuff2}\prep.

%%blk intro.opengap.head
\section{Why the comparison is still an open question}
\label{sec:opengap}

%%blk intro.opengap.01
The design space is known and the complexities are computed. A reader
could therefore assume that the comparison of the protocols is a solved
problem. It is not, for three reasons that the literature itself
documents.

%%blk intro.opengap.02
\paragraph{First: the published numbers do not transfer.}
Figure~\ref{fig:scaling} shows this in the hardest possible case for the
objection: the same work, the same unified implementation, the same
metric, and the same range of $n$. Only the deployment topology changes.
The scaling cost of PBFT moves from 65\% to 84\%, that of PoE from 39\%
to 51\%, and that of FTB from 32\% to 36\% \cite{amiri2024bedrock}. The
same work also shows a reversal of ranking with no change of environment
at all. SBFT gives 12\% \emph{higher} throughput than HotStuff at $n=4$,
and 8\% \emph{lower} at $n=100$, in the same experiment
\cite{amiri2024bedrock}.

%%blk intro.opengap.03
When the implementation changes as well, the differences become orders of
magnitude. In a simulation of a planetary network with $n=256$ nodes,
Kauri gives about twenty times the throughput of HotStuff-SECP256K1, that
is 4,917 operations per second against 246. Here each protocol has a
different implementation, and the values are simulated and not measured
\cite{berger2023simulation}.

%%blk intro.opengap.fig.scaling
\dgfigure{fig:scaling}%
{The same experiment in two deployment topologies: the percentage
reduction of throughput as the nodes rise from 4 to 100.}%
{The same experiment in two deployment topologies: the percentage
reduction of throughput as the nodes rise from 4 to 100. All the numbers
come from the same work and the same unified implementation of the
protocols \cite{amiri2024bedrock}. \textbf{(a)}~Deployment in one
datacenter: Amazon EC2, \texttt{c4.2xlarge} instances (8 vCPU,
15\,GB RAM, Intel Xeon E5-2666~v3 at 3.50\,GHz), one node per virtual
machine, batch size 400, payload 128/128\,bytes, clients in a closed loop
\cite{schroeder2006openclosed}\unver, mean of five runs, no faults. Three
clarifications for the figure. FaB needs $5f{+}1$, so it starts at $n=6$.
The work \cite{amiri2024bedrock} implements HotStuff in chained form,
with pipelining. FTB and FLB are new protocols that the same work
introduces, and not established ones. \textbf{(b)}~The three protocols
for which the same work also gives a geo-distributed number (App.~D, same
implementation framework, same metric, same batch of 400, no faults).
\emph{Two reservations that the work itself must state.} In the
geo-distributed set, \cite{amiri2024bedrock} raises the pipelining
stretch of Kauri and FTB from 4 to 6, and the range of FTB reaches $n=99$
instead of $n=100$. Neither of the two concerns PoE, so the clean
comparisons --- with no other change at all --- are two: \textbf{PBFT}
(65\,\% $\to$ 84\,\%) and \textbf{PoE} (39\,\% $\to$ 51\,\%). For all
three, however, the algorithm and the metric stay the same, and the
deployment topology is what changes.}%
{% fig-scaling — the same experiment, two deployment topologies. Two panels, so
% two dgbarchart environments stacked: the standard fixes every figure at
% \dgwidth, which leaves no room to set them side by side.
% All seven protocols are Byzantine, so the LAN series is dgBFT in both panels.
% Panel (b)'s second series is a second DEPLOYMENT of the same Byzantine
% protocols, so it takes a further-series colour (dgAux) — dgCFT would claim
% these are crash-fault protocols.
\begin{dgbarchart}[
  title={\fl{scaling.panel.lan}},
  ylabel={\fl{scaling.ylabel} (\%)},
  ymax=78,
  symbolic x coords={PBFT,FaB,SBFT,HotStuff,PoE,FTB,Kauri}, xtick=data,
  nodes near coords, every node near coord/.append style={dgtag, fill=dgPaper},
]
\addplot[dgBFT, fill=dgBFT] coordinates
  {(PBFT,65) (FaB,63) (SBFT,55) (HotStuff,45) (PoE,39) (FTB,32) (Kauri,31)};
\end{dgbarchart}

\par\vspace{5mm}

\begin{dgbarchart}[
  title={\fl{scaling.panel.geo}},
  ylabel={\fl{scaling.ylabel} (\%)},
  ymax=104,
  symbolic x coords={PBFT,PoE,FTB}, xtick=data,
  legend pos=north east,
  nodes near coords, every node near coord/.append style={dgtag, fill=dgPaper},
]
\addplot[dgBFT, fill=dgBFT,
         every node near coord/.append style={xshift=-2pt}]
  coordinates {(PBFT,65) (PoE,39) (FTB,32)};
  \addlegendentry{LAN}
\addplot[dgAux, fill=dgAux,
         every node near coord/.append style={xshift=2pt}]
  coordinates {(PBFT,84) (PoE,51) (FTB,36)};
  \addlegendentry{\fl{scaling.legend.geo}}
\end{dgbarchart}
}

%%blk intro.opengap.04
\paragraph{Second: the same comparison reverses, depending on what is
measured.}
The most instructive example is EPaxos. The original work claims optimal
commit latency in a wide-area network, uniform load distribution and
graceful degradation, with experiments on Amazon EC2
\cite{moraru2013epaxos}. The independent re-evaluation of Tollman et al.\
reaches three conclusions \cite{tollman2021epaxosrevisited}. It
\emph{confirms} the claim about \emph{median} latency. It shows, however,
a tail latency more than four times worse than Multi-Paxos. It also shows
that the share of operations that finish in one round trip fell to 65\%
in one of their experiments, under a load with high command conflict that
the original methodology did not produce. Matte et al.\ change the metric
again, from absolute throughput to efficiency of resources. They conclude
that in an environment of a fixed budget of resources, typical of the
real clouds, unoptimized Multi-Paxos reaches \emph{almost} twice the
throughput of EPaxos \cite{matte2021wasteful}. All three works are
correct. They simply do not measure the same thing. The consequence for
every comparative evaluation is that the choice of the metric is not a
technical detail but part of the conclusion, and it must be declared
before the measurement.

%%blk intro.opengap.05
\paragraph{Third: the difference of the protocol is not where the
measurement looks for it.}
Howard and Mortier analyse the two protocols that dominate the production
systems. They conclude that Paxos and Raft take \enquote{a very similar
approach to distributed consensus, differing only in their approach to
leader election} \cite{howard2020paxosraft}. If this holds, then a
measurement that compares two \emph{systems} which implement these
protocols measures, on the fault-free path, engineering of the
implementation and not an algorithmic difference. Such a measurement
therefore \emph{isolates} the implementation, which is exactly the right
experiment if that is the question. On the path of the leader change the
opposite holds. The whole algorithmic difference gathers there, and the
authors themselves make two separate statements. The first is a
structural observation: the election phase of Raft is \emph{lighter in
messages}, because it does not transfer entries of the log. The second is
an explicit prediction: they expect the election to be \emph{both slower
and of higher variance} in time \cite{howard2020paxosraft}. Two
reservations are obligatory. The comparison runs on a \emph{simplified}
statement of Paxos in the terminology of Raft, and the authors declare
that both algorithms they present are naive by design
\cite{howard2020paxosraft}.

%%blk intro.opengap.06
The gap, finally, has been stated explicitly, and it has been covered in
part. Ailijiang et al.\ find that the protocols of strong consistency
\enquote{show widely varying performance under different network,
workload, and deployment conditions, and no previous study offers a
comprehensive dissection and comparison of their performance}, and they
cover it for the crash-fault family \cite{ailijiang2019dissecting}. Amiri
et al.\ do the same for the Byzantine family \cite{amiri2024bedrock}.
Alqahtani and Demirbas measure protocols of both families in the same
framework \cite{alqahtani2021bottlenecks}.

%%blk intro.opengap.07
What stays open is therefore not the unified measurement in itself. It is
the measurement of the \emph{production implementations} --- the ones
that are really deployed --- with one instrument, one load model and
fault scenarios. All three works above measure research prototypes, or
unified reimplementations inside one framework, and not the systems that
run in production. As Figure~\ref{fig:scaling} shows, the
reimplementation is exactly the variable that can move the result.

%%blk intro.aims.head
\section{Aim, research questions and limits of the thesis}
\label{sec:aims}

%%blk intro.aims.01
The approved thesis proposal (see Appendix~A) gives the thesis three
aims. The first is the study of the evolution of the consensus
protocols. The second is the theoretical and architectural analysis of
the main ones. The third is the comparison of the performance of selected
implementations in representative scenarios. The comparison has four
criteria: consensus latency, throughput, fault tolerance and scaling. The
research questions below make these aims specific.

%%blk intro.aims.rq
\begin{description}[leftmargin=2.1em,style=nextline,font=\normalfont\bfseries]
  \item[RQ1 --- Evolution.] Which specific trade-offs explain each move
  inside the design space, and what price does each one pay? The answer
  is analytic and comes from the primary sources.
  \item[RQ2 --- Common structure.] Which invariants do all the protocols
  share, so that one single scheme can describe them instead of a list of
  independent algorithms?
  \item[RQ3 --- Comparability.] Under which conditions is a direct
  comparison of two implementations legitimate, and when does the
  measurement capture engineering of the implementation instead of a
  property of the protocol?
  \item[RQ4 --- Empirical behaviour.] How do selected production
  implementations behave in a controlled environment of a small cluster,
  in latency, throughput and fault tolerance? This thesis reports its own
  measurements for this in the evaluation chapter.
  \item[RQ5 --- Scaling.] How does the behaviour change as the quorum
  grows? The thesis approaches the question on two levels, and it states
  them separately. The experimental level covers only small cluster
  sizes, where the thesis tests the cost of the arithmetic of
  Figure~\ref{fig:quorum}. The literature level covers the large sizes,
  where the asymptotic behaviour appears. This thesis reports its own
  measurements for the experimental level in the evaluation chapter.
\end{description}

%%blk intro.aims.02
The thesis states its limits here and not in the conclusions. First, the
thesis evaluates \emph{implementations} and not protocols. The object
that can be measured is always a system with a specific version and a
specific configuration. Every claim about the underlying protocol must
therefore pass through RQ3.

%%blk intro.aims.03
Second, the available budget for infrastructure limits the range of
cluster sizes that an experiment can study. That range is small against
the scales where the interesting asymptotic behaviour appears. This is
the reason for the split of RQ5 into two levels.

%%blk intro.aims.04
Third, the thesis can cover a protocol \emph{analytically}, with no
measurements of our own. Where this happens, the chapter states it at its
start. Any numbers there come from third publications that the text names
explicitly, and they never stand in the same table as our own
measurements. The full framework of validity, the methodological criteria
and the threats to validity belong to the methodology chapter. This
chapter does not anticipate them.

%%blk intro.structure.head
\section{Contribution and structure of the thesis}
\label{sec:structure}

%%blk intro.structure.01
The thesis aims at four contributions.

%%blk intro.structure.list
\begin{enumerate}[label=\textbf{(\alph*)},leftmargin=2.4em]
  \item An account of the evolution of the protocols in time, with an
  explicit rule of order. The account attributes each transition to a
  named trade-off.
  \item One unified structural description. The same reading scheme
  covers the protocols of both fault models.
  \item A comparative evaluation of selected production implementations
  in a controlled environment, with pre-registered predictions. The
  thesis also states explicitly what the data \emph{cannot} support. This
  thesis reports its own measurements for this in the evaluation chapter.
  \item A decision framework that maps requirements to families of
  protocols.
\end{enumerate}

%%blk intro.structure.02
The structure follows two paths in time, one for each fault model,
instead of one strictly unified path. The choice is a deliberate
departure from the strictly chronological presentation, and it is put
explicitly to the supervisor. With one single order, the sequence of
Figure~\ref{fig:timeline} forces the reader to change fault model five
times. The result is that quantities which do not compare stand side by
side.

%%blk intro.structure.03
Chapter~\ref{ch:background} establishes the theoretical background: fault
and synchrony models, SMR, quorums, and criteria of correctness. The
chapters of the first path cover the crash-fault family, and the chapters
of the second path cover the Byzantine family. Each protocol follows one
common template of five subsections:

%%blk intro.structure.template
\begin{enumerate}[label=\arabic*.,leftmargin=2.4em]
  \item the historical context and the primary source;
  \item the algorithm;
  \item the implementation that was evaluated, or the reason for its
  absence;
  \item the measurements;
  \item the reading of the results against the prediction that was
  declared in advance.
\end{enumerate}

%%blk intro.structure.04
The chapter on methodology and validity comes next. After it come the
aggregate comparison, the decision framework and, last, the conclusions
with the limitations and the future work.

%%blk intro.sources.head
\section{A note on the quality of the sources}
\label{sec:sources}

%%blk intro.sources.01
A large part of the literature on consensus protocols circulates as a
preprint, as a whitepaper, or as technical documentation of a project.
For that reason the thesis applies three rules in all its chapters.
First, each algorithm is described from its primary source. The surveys
and the unified platforms serve for unifying schemes and for
classifications. Where such a work carries a bound or a number of a third
party, which happens only by exception, the text names the primary source
as well, as in the caption of Figure~\ref{fig:quorum}.

%%blk intro.sources.02
Second, the sources that have not been refereed are declared as such at
the point where they are used. The marker \grey\ shows the grey
literature: whitepapers, KIP documents and web pages. The marker \prep\
shows the preprints and the extended abstracts without review. A third
marker, \unver, says something different from those two: it says that the
bibliographic data of the source itself could not be checked against the
article or against an index, and were taken from the reference list of
another work. The reader who wants to follow that citation needs to know
this before trying. Third,
every published number that the thesis reproduces carries the environment
that produced it: hardware, load model, deployment topology, cluster size
and number of repetitions. The captions of Figures~\ref{fig:quorum}
and~\ref{fig:scaling} do this. A number without its environment is not
evidence.

%%blk intro.sources.03
The same principle applies to the bibliographic data. In the bibliography
file each entry carries a marker for how it was verified: from the first
page of the article itself, from an authentic index through the DOI, or
not at all. In the last case the entry is declared as unverified. The
distinction is not formalism. Three cases where the marking prevented an
error are worth recording.

%%blk intro.sources.04
The first case is \cite{berger2023simulation}. The copy that was
collected is a preprint and would have been declared as such, while the
work has a refereed version at a conference. The \emph{verification}
marker describes the copy that we hold. The \prep\ and \grey\ markers
describe the \emph{work}. Here the two diverge, and only a check against
an index finds the divergence. The copy is a preprint, the work is not,
and therefore \cite{berger2023simulation} does \emph{not} carry \prep.

%%blk intro.sources.05
The second case is \cite{oki1988viewstamped}. Crossref records it with a
different wording of the title from the wording of the proceedings, and
the bibliography of \cite{lamport1998parttime} itself confirms the full
form. The third case is \cite{jalalzai2023fasthotstuff}. Its title is
\enquote{\dots A Fast and Robust BFT Protocol}, not \enquote{Resilient},
as some catalogues show it.

%%blk intro.sources.06
A fourth case shows the same thing at the level of characters. The corpus
holds \emph{two} copies of \cite{lamport1998parttime}, the version of the
ACM and the corrected version of the author. They print the same
information with different letters: \enquote{Received: January 1990;
accepted: March 1998} in the first, \enquote{Received January 1990;
Accepted March 1998} in the second. Whichever one is quoted verbatim will
be wrong for the other. For that reason the caption of
Figure~\ref{fig:timeline} gives the \emph{content}. When two copies
differ, a verbatim quotation must name the copy as well.
